\section{Introduction}
\label{sec:intro}

Vision--language models (VLMs) are commonly evaluated for fine-grained understanding with image--text matching (ITM) benchmarks. In \svo{}~\cite{hendricks2021probing}, a caption comes with a positive and a negative image in which the subject, the verb or the object differs from the caption. The answer of a model is a binary classification of the pair, which a matching head, a prompted generative VLM or a re-ranker gives by thresholding its score. The understanding of a model is scored by its ITM accuracy, the share of balanced pairs it answers correctly~\cite{hendricks2021probing,bugliarello2023measuring}. We call such a model, together with its threshold, a matching decision.

Matching decisions are also used inside retrieval systems. ALBEF, BLIP and BLIP-2 retrieve in two stages: the image and text embeddings of their encoders rank the whole collection, and the matching head then re-ranks the top candidates, 128 or 256 of them~\cite{li2021align,li2022blip,li2023blip2}. Text retrieval has the same design, with a cross-encoder as the second stage~\cite{nogueira2019passage}, and Qwen3-VL-Reranker is released for this stage~\cite{qwen3vlembedding2026}. Whoever builds such a system has to choose the model for this stage, and also a threshold if it is to filter candidates. The candidates that reach this stage are similar to the query and to each other, since they were selected for this similarity: in \cref{fig:example}, the images retrieved for ``The man sits on the sofa.'' all show a man and a sofa, and the incorrect ones differ in what the man does. Telling such candidates apart is the fine-grained understanding that ITM benchmarks test, so one would expect a matching decision with a higher ITM accuracy to re-rank them better. The retrieved pairs, however, are not the benchmark's pairs: many are correct but unlabelled, and the incorrect ones are wrong in ways no benchmark designer wrote down.

\begin{figure}[t]
\centering
\includegraphics[width=\linewidth]{fig_example}
\caption{The caption ``The man sits on the sofa.'': its benchmark pair and six of the ten images FLAVA retrieves for it, with the human label and each head's match probability ($\blacktriangleleft$: the image it puts first). Both heads answer the benchmark pair correctly, yet the FLAVA head puts a verb error first and the BLIP-2 head the positive from rank 9.}
\label{fig:example}
\end{figure}

To test the expectation, we use \svo{} in a second role, as a retrieval collection. The benchmark's annotations pair each caption with one positive image, so a Success@1 that counts only these annotated pairs as correct treats every other retrieved item as incorrect, including the many that match the query~\cite{chun2022eccv,parekh2021crisscrossed}. We therefore reuse the judged lists of a companion study, under review, that compares embedding models: the top-10 lists of five of these models for 100 caption and 100 image queries, every query--candidate pair labelled by human annotators as correct or incorrect.

We ask three questions. (RQ1) How well does each decision separate correct from incorrect retrieved pairs? (RQ2) Does re-ordering a top-10 list by it improve the judged Success@1? (RQ3) Which errors does it let through? Our experiments show that ITM accuracy orders the decisions correctly but predicts neither the size nor the direction of their effect on the lists. We release the stored decisions and the analysis code with the judged lists.\footnote{\repolink}

\paragraph{Related work.}
Matching benchmarks with handcrafted negatives cover many phenomena beyond verbs~\cite{parcalabescu2021valse,thrush2022winoground,ma2023crepe,yuksekgonul2023when}, and their negatives carry artefacts that a text-only model can exploit~\cite{hsieh2023sugarcrepe}. Two studies show that a gain on such pairs need not carry over: when BiVLC adds the text-to-image direction to SugarCrepe, the conclusions of earlier work change~\cite{miranda2024bivlc}, and fine-tuning CLIP on hard negatives raises its benchmark scores while making it reject hard positives~\cite{kamath2024hardpositive}. Retrieve-and-re-rank pipelines with a cross-attention re-ranker were studied for image--text retrieval~\cite{miech2021thinking,geigle2022retrieve}, and generative VLMs are now prompted or trained as the re-ranker~\cite{lin2025mmembed,liu2025lamra,qwen3vlembedding2026}. These re-rankers are evaluated on the annotated pairs of Flickr30K and COCO, where every unannotated candidate counts as wrong although many match~\cite{parekh2021crisscrossed,chun2022eccv}. To our knowledge, no previous work has asked whether the ITM accuracy of a matching decision predicts its effect on judged retrieved lists.
